\documentclass[12pt,a4paper]{article}
\usepackage[UTF8]{ctex}
\usepackage{amsmath,amssymb}
\usepackage{geometry}
\geometry{left=2.5cm, right=2.5cm, top=2.5cm, bottom=2.5cm}
\usepackage{enumitem}
\usepackage{booktabs}

\title{校园驿站快递配送---多无人机协同调度\\数学模型}
\author{第三小组}
\date{2026年6月}

\begin{document}
\maketitle

% ============================================================
\section{问题描述}
% ============================================================

中山大学深圳校区快递驿站有 $m$ 件快递需由 $n$ 架美团 M4 无人机配送至校内各用户配送点。每件快递有重量 $w_j$、服务时间窗 $[a_j, b_j]$；每架无人机有最大载重 $W_i$、巡航速度 $v_i$。

无人机可连续执行多趟配送，每趟从起飞到返回起降点最长 2 小时（含飞行、送货和返回时间），之后需换电池（5 分钟）才能继续。每趟从起降点出发，依次飞往各取货点取货、飞往各配送点送货，最后返回起降点。

需确定：启用哪些无人机、每架配送哪些快递、按什么顺序、何时出发返回，使总配送成本最低。

% ============================================================
\section{集合与参数}
% ============================================================

\begin{table}[h]
\centering
\caption{主要符号说明}
\begin{tabular}{ll}
\toprule
符号 & 含义 \\
\midrule
$I = \{1,2,\ldots,n\}$ & 无人机集合 \\
$J = \{1,2,\ldots,m\}$ & 快递任务集合 \\
$R_i$ & 无人机 $i$ 的候选调度方案集合 \\
$s$ & 快递驿站（起降点） \\
$c_j$ & 快递 $j$ 的用户配送点 \\
$p_j$ & 快递 $j$ 的取货点 \\
$w_j$ & 快递 $j$ 的重量 (kg) \\
$[a_j, b_j]$ & 快递 $j$ 的服务时间窗 \\
$d_{uv}$ & 节点 $u$ 到 $v$ 的实际飞行距离 (m)，含绕行系数 1.08 \\
$v_i$ & 无人机 $i$ 的巡航速度 (m/s) \\
$W_i$ & 无人机 $i$ 的最大载重 (kg) \\
$F$ & 启用一架无人机的固定成本 \\
$T_{\max}$ & 单趟最长待机时间 (h)，从起飞到返回起降点，取 2 h \\
$T_{\text{swap}}$ & 换电池时间 (min)，取 5 min \\
$t_{\text{del}}$ & 每件快递送货时间 (min)，取 0.5 min \\
$\alpha$ & 时间成本权重 \\
$\gamma$ & 超时惩罚权重 \\
\bottomrule
\end{tabular}
\end{table}

% ============================================================
\section{决策变量}
% ============================================================

决策变量均为 0--1 变量：

\begin{itemize}[leftmargin=2em]
  \item $x_i$：无人机 $i$ 是否被启用
  \item $y_{ir}$：无人机 $i$ 是否选择并执行候选调度方案 $r$
\end{itemize}

% ============================================================
\section{候选调度方案生成}
% ============================================================

一个候选方案 $r \in R_i$ 表示无人机 $i$ 的一次完整调度，可能包含多趟配送。每趟从起降点出发，依次飞往各取货点取货、飞往各配送点送货，最后返回起降点。

\subsection{单趟路线}

无人机 $i$ 的一趟配送路线为：
\[
s \to (p_{j_1} \to c_{j_1}) \to (p_{j_2} \to c_{j_2}) \to \cdots \to (p_{j_q} \to c_{j_q}) \to s
\]

该路线需满足：

\textbf{(C1) 载重约束：}
\[
\sum_{\ell=1}^{q} w_{j_\ell} \leq W_i
\]

\textbf{(C2) 时限约束：} 每件快递的送达时间在时间窗内
\[
a_j \leq t_j^{\text{arrive}} \leq b_j, \quad \forall j
\]
早于 $a_j$ 到达则等待，晚于 $b_j$ 到达则记录超时量。

\textbf{(C3) 待机时间约束：} 单趟从起飞到返回起降点的总时间不超过 $T_{\max}$。
\[
t_{\text{flight}} + t_{\text{deliver}} + t_{\text{return}} \leq T_{\max}
\]

\subsection{多趟调度方案}

若无人机执行多趟配送，各趟之间需考虑换电池时间：
\[
\text{累计飞行时间} \geq T_{\max} \implies \text{换电池} \; T_{\text{swap}} \text{，之后继续}
\]

% ============================================================
\section{模型的标准形式}
% ============================================================

问题的完整数学规划形式如下：

\[
\boxed{
\begin{aligned}
\min \quad & Z = \alpha \sum_{i \in I} \sum_{r \in R_i} t_{ir} y_{ir} + F \sum_{i \in I} x_i + \gamma \sum_{i \in I} \sum_{r \in R_i} L_{ir} y_{ir} \\[8pt]
\text{s.t.} \quad
& \sum_{i \in I} \sum_{r \in R_i} a_{jr} y_{ir} = 1, & \forall j \in J \quad \text{(任务覆盖)} \\[4pt]
& y_{ir} \leq x_i, & \forall i \in I, r \in R_i \quad \text{(启用关联)} \\[4pt]
& \sum_{j \in J} w_j a_{jr} \leq W_i, & \forall i \in I, r \in R_i \quad \text{(载重约束)} \\[4pt]
& t_{ir} \leq T_j, & \forall i \in I, r \in R_i, j \in r \quad \text{(时限约束)} \\[4pt]
& x_i, y_{ir} \in \{0, 1\}, & \forall i \in I, r \in R_i
\end{aligned}
}
\]

其中：
\begin{itemize}[leftmargin=2em]
  \item $t_{ir}$：无人机 $i$ 执行方案 $r$ 的总飞行时间，包含飞行、送货和返回时间
  \item $L_{ir}$：无人机 $i$ 执行方案 $r$ 的总超时量（分钟）
  \item $a_{jr}$：路线 $r$ 是否覆盖任务 $j$（0--1 参数）
\end{itemize}

% ============================================================
\section{模型的目标}
% ============================================================

模型目标为最小化总配送成本：
\[
\min Z = \underbrace{F \sum_{i \in I} x_i}_{\text{启用成本}} + \underbrace{\alpha \sum_{i \in I} \sum_{r \in R_i} t_{ir} y_{ir}}_{\text{飞行时间成本}} + \underbrace{\gamma \sum_{i \in I} \sum_{r \in R_i} L_{ir} y_{ir}}_{\text{超时惩罚}}
\]

\begin{itemize}[leftmargin=2em]
  \item $F \cdot x_i$：启用成本，每启用一架无人机支付固定成本 $F = 0.5$
  \item $\alpha \cdot t_{ir}$：飞行时间成本，权重 $\alpha = 1.0$
  \item $\gamma \cdot L_{ir}$：超时惩罚，权重 $\gamma = 2.0$
\end{itemize}

% ============================================================
\section{模型的约束}
% ============================================================

模型约束如下：

\begin{itemize}[leftmargin=2em]
  \item \textbf{任务覆盖：} $\displaystyle\sum_{i \in I} \sum_{r \in R_i} a_{jr} y_{ir} = 1, \; \forall j \in J$，每件快递恰好被一个方案覆盖。
  \item \textbf{启用关联：} $y_{ir} \leq x_i, \; \forall i \in I, r \in R_i$，无人机 $i$ 须启用才能执行路线。
  \item \textbf{载重约束：} $\displaystyle\sum_{j \in J} w_j a_{jr} \leq W_i, \; \forall i \in I, r \in R_i$，每趟初始装载不超过最大载重。
  \item \textbf{时限约束：} $t_{ir} \leq T_j, \; \forall i \in I, r \in R_i, j \in r$，每件快递须在时限内送达。
\end{itemize}

% ============================================================
\section{求解方法}
% ============================================================

采用两阶段求解：

\begin{enumerate}[leftmargin=2em]
  \item \textbf{第一阶段---候选方案生成：} 枚举所有可行的单趟路线（满足载重、时限约束），再按时间顺序贪心组合多趟方案（满足电池时间约束）。对每个候选方案，预计算覆盖向量 $a_{jr}$、总时间 $t_{ir}$。
  \item \textbf{第二阶段---整数规划求解：} 以候选方案为决策空间，建立 0--1 整数规划模型，使用分支定界法（intlinprog）求解全局最优解。
\end{enumerate}

大规模问题（快递件数 $> 30$）采用启发式算法：

\begin{itemize}[leftmargin=2em]
  \item \textbf{随机搜索（Random Search）：} 对每种无人机数量（1--$n$ 架），各随机生成 100 个分配方案，用统一评估函数检查可行性，取全局最优。简单高效，适合作为基准。
  \item \textbf{遗传算法（Genetic Algorithm）：} 种群大小 100，迭代 200 代。采用锦标赛选择（$k=5$）、均匀交叉（概率 0.85）、多点变异（概率 0.15）。精英保留策略（5\%），以随机搜索最优解作为种子个体。搜索范围广，适合跳出局部最优。
\end{itemize}

所有启发式算法共用统一的评估函数 \texttt{eval\_solution}，确保约束检查和成本计算的一致性。

% ============================================================
\section{统一评估函数}
% ============================================================

所有算法共用 \texttt{eval\_solution(assignment, data, cfg, dep)} 函数，硬性约束如下：

\begin{enumerate}[leftmargin=2em]
  \item \textbf{所有订单必须被分配：} assignment 向量中不能有 0
  \item \textbf{载重约束：} 每趟总重量 $\leq W_i = 2.5$ kg，超重则返回起降点开新趟
  \item \textbf{电池约束：} 每趟总飞行时间 $\leq T_{\max} = 2$ h，超时则返回起降点开新趟
  \item \textbf{时间窗约束：} 到达时间 $= \max(\text{订单就绪时间}, \text{无人机可用时间}) + \text{飞行时间}$
  \item \textbf{换电池时间：} 两趟之间间隔 $T_{\text{swap}} = 5$ min
\end{enumerate}

% ============================================================
\section{实验结果}
% ============================================================

\subsection{数据规模}

\begin{table}[h]
\centering
\caption{实验数据}
\begin{tabular}{lr}
\toprule
项目 & 数量 \\
\midrule
无人机数量 & 8 \\
快递订单数 & 120 \\
配送节点数 & 9（2 起降点 + 4 取货点 + 3 配送点） \\
候选路线数 & 256（route\_info.csv） \\
\bottomrule
\end{tabular}
\end{table}

\subsection{算法对比}

\begin{table}[h]
\centering
\caption{算法对比结果（120 订单，实际飞行距离）}
\begin{tabular}{lccccc}
\toprule
方法 & 总成本 & 启用数 & 超时(min) & 总用时(min) & 计算耗时(s) \\
\midrule
随机$\times$100/架 & 5.3 & 2 & 6.5 & 243.6 & 1.25 \\
遗传算法 & \textbf{4.8} & \textbf{2} & \textbf{0.0} & \textbf{229.9} & 20.67 \\
\bottomrule
\end{tabular}
\end{table}

遗传算法在种群多样性与交叉重组机制下，能够找到更优的配送方案，总成本较随机搜索降低 9.4\%，且无超时订单。

\subsection{300 订单扩展实验}

\begin{table}[h]
\centering
\caption{算法对比结果（300 订单）}
\begin{tabular}{lccccc}
\toprule
方法 & 总成本 & 启用数 & 超时(min) & 总用时(min) & 计算耗时(s) \\
\midrule
随机$\times$100/架 & 12.5 & 5 & 0.7 & 601.3 & 3.58 \\
遗传算法 & \textbf{12.0} & \textbf{5} & \textbf{0.0} & \textbf{567} & 59.20 \\
\bottomrule
\end{tabular}
\end{table}

300 订单规模下，遗传算法依然保持优势，总成本 12.0，启用 5 架无人机，零超时。

% ============================================================
\section{无人机参数}
% ============================================================

\begin{table}[h]
\centering
\caption{美团 M4 无人机参数}
\begin{tabular}{ll}
\toprule
参数 & 值 \\
\midrule
最大载重 $W_i$ & 2.5 kg \\
巡航速度 $v_i$ & 15 m/s（54 km/h） \\
单趟最长待机时间 $T_{\max}$ & 2 h \\
换电池时间 $T_{\text{swap}}$ & 5 min \\
每件送货时间 $t_{\text{del}}$ & 0.5 min \\
启用固定成本 $F$ & 0.5 \\
时间成本权重 $\alpha$ & 1.0 \\
超时惩罚权重 $\gamma$ & 2.0 \\
\bottomrule
\end{tabular}
\end{table}

% ============================================================
\section{距离计算}
% ============================================================

本项目使用实际飞行距离（含绕行系数 1.08），而非直线距离。距离矩阵来自 \texttt{distance\_matrix.csv}，基于中山大学深圳校区节点坐标计算。

\begin{table}[h]
\centering
\caption{节点间实际飞行距离示例（单位：m）}
\begin{tabular}{lcccc}
\toprule
起点 & 终点 & 实际距离 & 直线距离 & 绕行系数 \\
\midrule
DEP\_W & CANTEEN\_W & 302.3 & 280.0 & 1.08 \\
DEP\_E & DORM\_E & 290.6 & 269.0 & 1.08 \\
RESTAURANT\_X & LIB & 183.6 & 170.0 & 1.08 \\
\bottomrule
\end{tabular}
\end{table}

\end{document}
